\chapter*{Présentation}
\addcontentsline{toc}{chapter}{Présentation}
Ce livre propose cinq introductions à des idées mathématiques qui permettent de représenter, de traiter et d'apprendre à partir des données. Les exemples concrets servent de point de départ, tandis que des énoncés précis expliquent ce qu'un algorithme peut ou ne peut pas garantir.
\begin{itemize}
\item Le chapitre~\ref{chap-shannon} introduit le codage binaire et la borne entropique de Shannon pour la compression sans perte.
\item Le chapitre~\ref{chap-images} considère une image comme un tableau de nombres et développe la quantification, le filtrage, la détection de contours et les transformations de couleurs.
\item Le chapitre~\ref{chap-sparsity} relie l'échantillonnage et l'approximation parcimonieuse à la régularisation des problèmes inverses et à l'acquisition comprimée.
\item Le chapitre~\ref{chap-ot} développe le transport optimal, des affectations finies à la comparaison et à l'interpolation de distributions.
\item Le chapitre~\ref{chap-neural-networks} explique les réseaux de neurones discriminatifs et génératifs, leurs objectifs d'apprentissage et leur lien avec le transport.
\end{itemize}
Les deux premiers chapitres demandent des connaissances d'algèbre élémentaire et, pour la théorie de l'information, de logarithmes et de probabilités de base. Les suivants utilisent les vecteurs, les matrices, les dérivées et les distributions de probabilité. Les résultats les plus avancés sont énoncés avec leurs hypothèses ; les exemples qui les entourent peuvent se lire avant d'étudier ces résultats en détail.

Chaque chapitre peut être lu séparément. Une lecture suivie fait apparaître des liens utiles : l'entropie explique le codage, les transformées rendent les signaux compressibles, la régularisation sélectionne les reconstructions et le transport compare les distributions qu'un modèle génératif cherche à apprendre. Une notation commune et de courtes introductions de chapitres accompagnent ces deux parcours.

Les éditions française et anglaise contiennent les mêmes cinq chapitres. Les livres complets, les chapitres séparés et les sources sont disponibles à l'adresse
\begin{center}\url{https://mathematical-tours.github.io/book-basics/}\end{center}
Pour un traitement plus approfondi, on pourra consulter \emph{Mathematical Foundations of Data Sciences}, disponible sur \url{https://mathematical-tours.github.io/book/}.

\paragraph*{Remerciements.}
Je remercie Marie-Noëlle Peyré, Gwenn Guichaoua, François Béguin, Gérard Grancher, Aurélien Djament, François Sauvageot, Charles Dossal, Jalal Fadili, Samuel Vaiter, Stéphane Seuret, Vincent Beck, Sébastien Racanière, Vincent Barra et le rapporteur anonyme de l'article original sur la parcimonie pour leurs relectures attentives et leurs corrections. La photographie de la fleur est de Maïtine Bergounioux.

\medskip
\noindent Gabriel Peyré\\
CNRS \& DMA, École normale supérieure, PSL
